\section{Related Work}

\subsection{Open Government Data und Bürgerbeteiligung}

Sayogo et al.\ \cite{open_gov} analysieren in einem Benchmarking-Framework
35 nationale Open-Data-Portale und identifizieren ein strukturelles Problem:
Daten werden häufig veröffentlicht, weil sie technisch einfach bereitzustellen sind –
nicht weil sie für die Bevölkerung besonders relevant wären.
Rund 66\,\% der Portale bieten Daten in maschinenlesbaren Formaten an,
doch eine nutzerorientierte visuelle Aufbereitung fehlt fast überall.
Diese Erkenntnis bildete die Grundlage unserer Designentscheidung,
Open Data nicht als Rohdaten, sondern als räumlich verortete,
kontextualisierte Information anzuzeigen.

\subsection{Dashboards und Visualisierung für Bürger:innen}

Simonofski et al.\ \cite{simonofski2022} evaluieren in einer Studie mit
108 Teilnehmenden 16 Designprinzipien für Open-Government-Data-Dashboards.
Ihr zentrales Ergebnis: Gut gestaltete Dashboards übertreffen isolierte
Einzelvisualisierungen deutlich – in Verständlichkeit, Usability und
wahrgenommener Datenqualität.
Besonders betonen sie die Notwendigkeit, \emph{heterogene Nutzergruppen}
explizit zu berücksichtigen: von technikaffinen Expert:innen bis zu
Gelegenheitsnutzer:innen ohne Datenkenntnisse.
Diese Erkenntnis spiegelt sich direkt in unseren Personas (Benjamin, Leah,
Ralph, Luca) und in der gestuften Informationsarchitektur von MapApp wider:
Die Karte bietet einen sofortigen visuellen Überblick;
Details sind in einem Bottom Sheet verborgen und werden erst bei
gezieltem Interesse angezeigt.

\subsection{Einprägsame Visualisierungen}

Borkin et al.\ \cite{Visualization} untersuchen in einer Studie
mit $n = 261$ Teilnehmenden, welche Eigenschaften Visualisierungen
einprägsam machen.
Zwei Ergebnisse sind für MapApp besonders relevant:
Erstens bleiben \emph{farbige} Visualisierungen mit erkennbaren Symbolen
signifikant besser in Erinnerung als abstrakte Diagramme.
Zweitens erhöhen \emph{ungewöhnlichere} Darstellungsformen die Merkfähigkeit.
Diese Erkenntnisse rechtfertigen unser farbkodiertes Pin-System:
Statt einer einheitlichen Darstellung aller Datenpunkte werden Pins nach
Kategorie eingefärbt (Orange~=~Kultur, Blau~=~Verkehr, Grün~=~Umwelt,
Lila~=~Infrastruktur, Grau~=~Basis), was die Orientierung auf der Karte
visuell einprägsam macht.

\subsection{Gestaltungsprinzipien für Datenvisualisierungen}

Rougier et al.\ \cite{simple_rules} formulieren zehn praktische Regeln
für gute Datenvisualisierungen.
Für MapApp besonders relevant sind Regel~3 (\textit{Beachte das Medium}):
Eine mobile App erfordert kompakte, auf kleinen Bildschirmen gut lesbare
Darstellungen – daher verzichten wir auf eingebettete Diagramme und
setzen stattdessen auf die Karte als primäres Visualisierungsmedium.
Regel~8 (\textit{Vermeide Chartjunk}) rechtfertigt das bewusste Verbergen
sekundärer Informationen im Bottom Sheet und im Swipe-to-reveal-Muster.

\subsection{Verwandte Anwendungen und Konkurrenzprodukte}

Unsere Konkurrenzanalyse (M1) zeigt, dass webbasierte Plattformen wie
\textit{Winterthur in Zahlen} und \textit{visualize.admin.ch} zwar
Kartenansichten bieten, aber nicht für mobile, kontextbezogene Nutzung
ausgelegt sind.
\textit{Eurostat} bietet die visuell ansprechendsten Darstellungen,
ist aber in seiner Individualiserbarkeit eingeschränkt.
Der \textit{OECD Data Explorer} leidet unter schwieriger und unflüssiger
Bedienung.
Die offizielle \textit{Wien Karte} \cite{wienkarte} der Stadt Wien deckt
zwar Transportdaten ab, bietet aber keine allgemeine Open-Data-Exploration
oder Datensatzverwaltung.
MapApp füllt diese Lücke: eine mobile, kartenbasierte Anwendung
speziell für Wiens Open-Data-Ökosystem.
